\chapter{Matplotlib による可視化}
\label{chap:matplotlib}

Matplotlib\index{matplotlib@Matplotlib|textbf} は、Python でグラフや画像を描画するライブラリである。本章では、データの性質に応じた図の選択と、オブジェクト指向 API による構成要素の制御を扱う。前半では基本的な作図と図種の選択、後半では Figure、Axes、Axis、Artist の関係、複数パネル、スタイル、保存方法を説明する。

ヒストグラム\index{ひすとぐらむ@ヒストグラム}や散布図には、計算結果の検証に用いる診断図としての役割がある。平均値や標準偏差だけでは識別できない長い裾、二峰性、外れ値などを、分布の形として確認できるためである。したがって、解析途中の検証と最終結果の提示の双方に可視化を組み込む。

\section{線図・散布図・ヒストグラムの基本}

\subsection{最小限のヒストグラム}

最初に、乱数生成器のシードを固定して正規乱数を 2500 個生成し、40 個のビンへ数え上げて PDF に保存する。入力は \code{values}、作図処理は \code{plt.hist()}、出力先は \code{histogram.pdf} である。

\begin{lstlisting}[language=Python]
import numpy as np
import matplotlib.pyplot as plt

rng = np.random.default_rng(7)
values = rng.normal(loc=0.0, scale=1.2, size=2500)

plt.hist(values, bins=40)
plt.xlabel("value")
plt.ylabel("count")
plt.savefig("histogram.pdf")
\end{lstlisting}

線、文字、マーカーを主体とする図は、ベクタ形式の PDF で保存すると拡大時にも輪郭を保てる。PDF は LaTeX 文書へ直接挿入できる。図~\ref{fig:basic-histogram-example} は、この最小例の出力である。

\code{plt.show()} は対話的な確認に、\code{savefig()} は出力条件を指定したファイル保存に用いる。保存先、ファイル形式、解像度、同名ファイルを上書きするかどうかを作図コード内で明示すると、図の生成条件を再現できる。

\begin{figure}[tbp]
\centering
\includegraphics[width=0.72\linewidth]{basic_histogram_example.pdf}
\caption{入力配列 \code{values} を 40 ビンへ数え上げたヒストグラム。横軸は値、縦軸は各ビンに含まれる要素数を表す。}
\label{fig:basic-histogram-example}
\end{figure}

\subsection{軸ラベル、凡例、目盛り}

図には、各軸が表す量と単位、系列を識別する凡例、比較に必要な目盛りを記す。これらが欠けると、数値の物理的意味や系列間の対応を図だけから判定できない。

\begin{lstlisting}[language=Python]
import numpy as np
import matplotlib.pyplot as plt

x = np.linspace(0.0, 10.0, 60)
y = np.exp(-x / 4.0) * (1.0 + 0.08 * np.sin(2.5 * x))

fig, ax = plt.subplots(figsize=(6, 4))
ax.plot(x, y, marker="o", label="sample A")
ax.set_xlabel(r"$t\,[\mathrm{s}]$")
ax.set_ylabel("signal")
ax.set_title("Labeled line example")
ax.set_xticks([0, 2, 4, 6, 8, 10])
ax.legend(loc="upper right")
fig.savefig("labeled_plot_example.pdf")
\end{lstlisting}

図~\ref{fig:labeled-plot-example} には、横軸の量と単位、縦軸の量、系列名、比較に必要な目盛りが示されている。配色を調整する前に、これらの情報が図だけから読み取れるかを確認する。

\begin{figure}[tbp]
\centering
\includegraphics[width=0.72\linewidth]{labeled_plot_example.pdf}
\caption{ラベル、凡例、目盛りを付けた線図。横軸は秒単位の時刻、縦軸は信号値、凡例は系列名を表す。}
\label{fig:labeled-plot-example}
\end{figure}

\subsection{plot 関数へ渡すデータ}

Matplotlib の多くの作図関数は、NumPy 配列や Python リストのような「配列に変換できるもの」を受け取る。加えて、辞書や DataFrame のように列名で参照できるオブジェクトなら、\code{data=} を使って列名指定で描くこともできる。

\begin{lstlisting}[language=Python]
import numpy as np
import matplotlib.pyplot as plt

times = [0, 1, 2, 3]
values = [1.0, 2.0, 1.5, 3.0]

fig, ax = plt.subplots(figsize=(6, 4))
x = np.asarray(times)
y = np.asarray(values)
ax.plot(x, y)

table = {"time": x, "value": y}
ax.plot("time", "value", data=table)
\end{lstlisting}

\code{np.asarray} は、入力を NumPy 配列として扱うことを明示する。\code{data=} を用いる場合は列名を文字列で指定するため、存在しない列名を渡したときのエラーも確認する。独立した配列を渡す場合は、\code{x} と \code{y} を明示する。

\subsection{線・点・参照線・誤差棒の描画関数}

線図には \code{plot()}\index{matplotlib@Matplotlib!plot@\code{plot()}|textbf}、散布図には \code{scatter()}\index{matplotlib@Matplotlib!scatter@\code{scatter()}|textbf}、鉛直および水平の参照線には \code{axvline()}\index{matplotlib@Matplotlib!axvline@\code{axvline()}} と \code{axhline()}\index{matplotlib@Matplotlib!axhline@\code{axhline()}}、誤差付きの点には \code{errorbar()}\index{matplotlib@Matplotlib!errorbar@\code{errorbar()}|textbf} を用いる。次の例は、これらの関数と出力図の対応を示す。

\begin{lstlisting}[language=Python]
import numpy as np
import matplotlib.pyplot as plt

rng = np.random.default_rng(1234)
t = np.linspace(0.0, 10.0, 80)
y = np.sin(t) * np.exp(-t / 12.0)
y_noise = y + 0.12 * rng.normal(size=t.size)
y_err = 0.10 + 0.03 * np.cos(t / 2.0) ** 2

fig, axes = plt.subplots(2, 2, figsize=(10, 7))

axes[0, 0].plot(t, y, marker="o")
axes[0, 1].scatter(t, y_noise, c=t, cmap="viridis")
axes[1, 0].plot(t, y_noise)
axes[1, 0].axvline(3.0, linestyle="--")
axes[1, 0].axhline(0.0, linestyle=":")
axes[1, 1].errorbar(t[::4], y_noise[::4], yerr=y_err[::4], fmt="o")

fig.savefig("basic_plot_types_example.pdf")
\end{lstlisting}

図~\ref{fig:basic-plot-types-example} は、主要な描画関数の出力を並べたものである。線図は順序を持つ系列や関数形、散布図は 2 変量の対応とばらつき、参照線は閾値や基準値、エラーバーは測定値の不確かさを表す。

\begin{figure}[tbp]
\centering
\includegraphics[width=0.94\linewidth]{basic_plot_types_example.pdf}
\caption{主要な描画関数の出力。左上は \code{plot()}、右上は \code{scatter()}、左下は \code{axvline()} と \code{axhline()}、右下は \code{errorbar()} の例である。}
\label{fig:basic-plot-types-example}
\end{figure}

\code{plot()} は既定で点を線で結ぶため、順序を持つ系列または連続関数の表示に用いる。\code{scatter()} は点どうしを結ばず、個々のイベントや標本の 2 変量分布を表示する。\code{axvline()} と \code{axhline()} は、閾値や基準値を対応する軸上へ示す。

\subsection{ヒートマップとカラーバー}

格子上で定義された 2 次元配列は \code{imshow()}\index{matplotlib@Matplotlib!imshow@\code{imshow()}|textbf} で色に対応させる。これに対して、点群を格子ごとに数え上げる場合は \code{hist2d()}\index{matplotlib@Matplotlib!hist2d@\code{hist2d()}|textbf} を用いる。いずれの場合も、色が表す量と尺度をカラーバーに記す。

\begin{lstlisting}[language=Python]
import numpy as np
import matplotlib.pyplot as plt
from matplotlib.colors import LogNorm

x = np.linspace(-3.0, 3.0, 160)
y = np.linspace(-2.5, 2.5, 140)
xx, yy = np.meshgrid(x, y)
image = np.exp(-(xx**2 + 1.6 * yy**2)) + 0.02
rng = np.random.default_rng(1234)
sample_x = rng.normal(loc=0.0, scale=1.0, size=80000)
sample_y = 0.7 * sample_x + rng.normal(loc=0.0, scale=0.8, size=80000)

fig, axes = plt.subplots(1, 2, figsize=(10, 4))
im = axes[0].imshow(image, origin="lower", cmap="viridis", aspect="auto")
fig.colorbar(im, ax=axes[0], label="intensity")

h = axes[1].hist2d(sample_x, sample_y, bins=90, norm=LogNorm(), cmap="magma")
fig.colorbar(h[3], ax=axes[1], label="counts")
\end{lstlisting}

図~\ref{fig:heatmap-colorbar-example} の左は \code{imshow()} によるヒートマップ、右は \code{hist2d()} と \code{LogNorm} を使った 2 次元ヒストグラムである。\code{cmap=} は値から色への対応\index{からーまっぷ@カラーマップ}を定め、\code{fig.colorbar(...)} はその対応尺度と量名を図中に示す。

\begin{figure}[tbp]
\centering
\includegraphics[width=0.94\linewidth]{heatmap_colorbar_example.pdf}
\caption{ヒートマップとカラーバーの例。左は \code{imshow()}、右は \code{hist2d()} と \code{LogNorm} を使った例であり、どちらもカラーバーによって色の意味を示している。}
\label{fig:heatmap-colorbar-example}
\end{figure}

\code{imshow()} は格子上に並んだ値を描くため、シミュレーションで得た場や画像データをそのまま表示できる。一方、\code{hist2d()} は点群を格子ごとに数え上げる。点数が数桁にわたる場合は、\code{LogNorm()} による対数的な色割り当てによって、高密度領域と低密度領域を同じ図に表示できる。

\subsection{全天図と投影: Hammer 図法と Mollweide 図法}

経度と緯度のような球面上のデータを平面へ描くには、投影法\index{とうえいほう@投影法}が必要になる。Matplotlib には、地図ライブラリなしでも使える全天図向けの投影として、Hammer\index{hammer@Hammer projection} 図法や Mollweide\index{mollweide@Mollweide projection} 図法が用意されている。

\begin{lstlisting}[language=Python]
import numpy as np
import matplotlib.pyplot as plt

lon_deg = np.linspace(-180.0, 180.0, 241)
lat_deg = 25.0 * np.sin(np.deg2rad(2.0 * lon_deg))
lon = np.deg2rad(lon_deg)
lat = np.deg2rad(lat_deg)

fig = plt.figure(figsize=(10, 4))
ax1 = fig.add_subplot(1, 2, 1, projection="hammer")
ax2 = fig.add_subplot(1, 2, 2, projection="mollweide")

ax1.scatter(lon, lat)
ax2.scatter(lon, lat)
\end{lstlisting}

図~\ref{fig:projection-examples} は、同じ点群を Hammer 図法と Mollweide 図法で投影した例である。どちらも経度の全範囲と両極を有限の領域に写すため、全天イベント分布や銀河座標図に利用できる。\code{projection="hammer"} と \code{projection="mollweide"} の座標値にはラジアンを渡す。

\begin{figure}[tbp]
\centering
\includegraphics[width=0.94\linewidth]{projection_examples.pdf}
\caption{Hammer 図法と Mollweide 図法の例。どちらも経度の全範囲と両極を一つの楕円領域内に投影する。}
\label{fig:projection-examples}
\end{figure}

国境、海岸線、地理座標系の変換を必要とする地図には Cartopy を用いる。球面座標上の点分布だけを表示する場合は、Matplotlib の標準投影を利用できる。入力座標の単位、経度の範囲、周期境界における処理を作図前に確認する。 Cartopy は追加のライブラリであり、導入方法は公式資料~\cite{cartopy_docs} を参照する。海岸線などの地図データを初めて使う際には取得が必要になる場合がある。


\section{分布・相関・密度に応じた図の選択}

2 変数の関係を調べるときに、データ量を確認せず散布図を描くだけでは、点の重なりによって密度構造を見落とす場合がある。この節の対象は、データ点数と確認したい性質に応じた散布図、透過表示、2 次元ヒストグラムの選択である。

\subsection{度数と確率密度の区別}
\index{ひすとぐらむ@ヒストグラム!せいきか@正規化|textbf}\index{かくりつみつど@確率密度|textbf}
\code{hist()} の既定の縦軸はビン内の度数である。\code{density=True} は、幅 \(\Delta x_i\) のビンの度数 \(n_i\) を \(n_i/(N\Delta x_i)\) に正規化する。ここで \(N\) は表示範囲内で数えた総度数であり、全ビンの面積の和が 1 になる。棒の高さの和が 1 になるとは限らない。

不等幅ビンや対数間隔のビンでは、広いビンほど多くの点が入り得る。分布の形を比較する場合は、度数、単位幅当たりの度数、確率密度のどれを示すかを軸に明記する。\code{density=True} は横軸の値そのものの幅で割る指定であり、対数軸にしただけで対数値の確率密度にはならない。確率密度関数を重ねるときも、正規化する範囲と単位をそろえる。

\subsection{点の重なりに応じた表示方法の選択}

\subsubsection{散布図による点の重なりの確認}

2 変数の対応を点ごとに示す基本的な方法が散布図（\code{scatter}）である。ただし、データ点が多い場合は同じ領域に点が重なり、局所的な密度差を判別できなくなる。

\begin{lstlisting}[language=Python]
import numpy as np
import matplotlib.pyplot as plt

# 約 60 万点を不透明な点として描画する
rng = np.random.default_rng(19)
x = rng.normal(loc=0.0, scale=1.0, size=600_000)
y = 0.65 * x + rng.normal(loc=0.0, scale=0.55, size=x.size)

fig, ax = plt.subplots(figsize=(6, 5))
ax.scatter(x, y, s=10, c='black')
fig.savefig("scatter_bad.png", dpi=150)
\end{lstlisting}

\begin{figure}[htbp]
\centering
\includegraphics[width=0.6\linewidth]{scatter_bad.png}
\caption{約 60 万点を不透明な点で描いた散布図。点の重なりによって、中心部の密度差を判別できない。}
\label{fig:scatter-bad}
\end{figure}

図~\ref{fig:scatter-bad} からは点が存在する範囲は分かるが、各領域の点数は読み取れない。領域間の密度は、点数をビンごとに集計する 2 次元ヒストグラムで比較する。

\subsubsection{2 次元ヒストグラムによる密度の計数}

2 次元ヒストグラム（\code{hist2d}）は、領域ごとの点数を色で表すため、点の重なりがある場合にも密度を定量的に示せる。

\begin{lstlisting}[language=Python]
fig, ax = plt.subplots(figsize=(6, 5))
h_lin = ax.hist2d(x, y, bins=50, cmap='viridis')
cbar = fig.colorbar(h_lin[3], ax=ax, label="Counts (Linear)")
fig.savefig("hist2d_linear_bad.pdf")
\end{lstlisting}

\begin{figure}[htbp]
\centering
\includegraphics[width=0.6\linewidth]{hist2d_linear_bad.pdf}
\caption{線形カラースケールを用いた 2 次元ヒストグラム。中心部の高密度領域に色範囲の大半が割り当てられ、低密度領域の差を識別しにくい。}
\label{fig:hist2d-linear}
\end{figure}

図~\ref{fig:hist2d-linear} では中心部の密度を読み取れるが、中心部と周辺部の度数が大きく異なるため、周辺の少数カウントの領域には近い色が割り当てられる。この線形カラースケールでは、周辺部の密度差や裾の広がりを識別しにくい。

\subsubsection{\code{LogNorm} による広い密度範囲の表示}

物理イベントの頻度が 1、10、100、1000 のように数桁にわたる場合には、線形カラースケールだけでは低密度領域の差を表しにくい。\code{LogNorm} は、正のカウントを対数尺度で色へ割り当てる。\index{matplotlib@Matplotlib!LogNorm@\code{LogNorm}|textbf}

\begin{lstlisting}[language=Python]
import matplotlib.colors as mcolors

fig, ax = plt.subplots(figsize=(6, 5))
# カウントから色への割り当てを対数尺度にする
h_log = ax.hist2d(x, y, bins=50, norm=mcolors.LogNorm(), cmap='viridis')
cbar = fig.colorbar(h_log[3], ax=ax, label="Counts (Log Scale)")
fig.savefig("hist2d_lognorm_good.pdf")
\end{lstlisting}

\begin{figure}[htbp]
\centering
\includegraphics[width=0.6\linewidth]{hist2d_lognorm_good.pdf}
\caption{\code{LogNorm} を用いた 2 次元ヒストグラム。中心の高密度領域と周囲の低密度領域を同じ色尺度で表す。}
\label{fig:hist2d-lognorm}
\end{figure}

図~\ref{fig:hist2d-lognorm}では、単一のピークの周囲に広がる低密度の広がりも判別できる。線形尺度では近い色に対応していた低密度領域が、対数尺度では異なる色に割り当てられるためである。

\subsection{\code{subplots()} による複数パネルの配置}

関連する分布を同じ尺度で比較する場合は、1 枚のキャンバス（\code{Figure}\index{matplotlib@Matplotlib!Figure@\code{Figure}|textbf}）の中に複数のパネル（\code{Axes}\index{matplotlib@Matplotlib!Axes@\code{Axes}|textbf}）を配置する。

\begin{lstlisting}[language=Python]
import numpy as np
import matplotlib.pyplot as plt

rng = np.random.default_rng(23)
data_a = rng.normal(loc=0.0, scale=1.0, size=2000)
data_b = rng.normal(loc=0.4, scale=1.3, size=2000)

# 1 行 2 列（左右に 2 つ）のパネルを作成する。
# fig は Figure 全体、axes は二つの Axes を持つ配列である。
fig, axes = plt.subplots(1, 2, figsize=(10, 4))

# 左のパネルには axes[0] を使って描画する
axes[0].hist(data_a, bins=30)
axes[0].set_title("Data A")

# 右のパネルには axes[1] を使って描画する
axes[1].hist(data_b, bins=30)
axes[1].set_title("Data B")

# 複数のパネルを描くと軸の数字やラベルが隣の図と重なりやすいため、
# 描画後に余白を調整する。
fig.tight_layout()
fig.savefig("histogram_pair_example.pdf")
\end{lstlisting}

複数パネルの図では、\code{axes[0]} のように描画先の \code{Axes} を明示する。これにより、各処理がどのパネルに作用するかをコード上で追跡できる。\code{Figure} と \code{Axes} の役割は、第~\ref{sec:matplotlib-object}~節以降で詳しく説明する。

\subsection{単純な横並び以外: \code{subplot\_mosaic} と \code{GridSpec}}

上段の主結果に広い領域を割り当て、下段に補助図を 2 枚置くなど、パネルごとに面積を変える場合は、\code{subplot\_mosaic()}\index{subplot_mosaic@subplot\_mosaic} または \code{GridSpec}\index{gridspec@GridSpec} を用いる。

\begin{lstlisting}[language=Python]
import numpy as np
import matplotlib.pyplot as plt

values = np.random.default_rng(42).normal(size=2000)
x = np.linspace(0.0, 10.0, 150)
y_top = np.exp(-x / 5.0) * np.cos(2.0 * x)
x_sc = np.linspace(0.0, 1.0, 25)
y_sc = x_sc**2 + 0.05 * np.sin(20.0 * x_sc)

layout = [["A", "A"], ["B", "C"]]
fig, axes = plt.subplot_mosaic(layout, figsize=(10, 6))

axes["A"].plot(x, y_top)
axes["B"].hist(values, bins=40)
axes["C"].scatter(x_sc, y_sc)
\end{lstlisting}

図~\ref{fig:subplot-mosaic-example} のように、\code{"A"} を 2 回書くと、上段の 2 列を一つの Axes が占める。主図と補助図に異なる面積を割り当てる場合に利用できる。

\begin{figure}[tbp]
\centering
\includegraphics[width=0.92\linewidth]{subplot_mosaic_example.pdf}
\caption{\code{subplot\_mosaic()} を使った複合レイアウトの例。上段は 2 列をまたぐ主図、下段は補助図 2 枚であり、単純な \code{plt.subplots(1, 2)} では作れない配置になっている。}
\label{fig:subplot-mosaic-example}
\end{figure}

\code{subplot\_mosaic()} は文字列の名前で Axes を参照できる。\code{GridSpec} は、列幅、行高、結合するセルを数値で指定する。Axes の参照方法と格子比率の制御範囲に応じて選択する。

\code{LogNorm} は 0 以下の値を対数へ写せないため、空のビンなどは通常マスクされる。背景と同じ色で表示された領域を、微小な正の度数と混同しないようにする。また、\code{imshow()} の既定の座標は画像の画素位置である。物理的な座標で表示する場合は、\code{origin} と \code{extent=(xmin, xmax, ymin, ymax)} で上下の向きと画像の外縁を明示する。

\subsection{凡例とレイアウト}

複数の系列やフィット曲線を同じ図へ重ねる場合は、各系列を識別する凡例が必要である。凡例の位置、枠の有無、項目の順序は、データとの重なりと系列間の対応を考慮して定める。

次の例では、同じ物理量を条件 A と B で測定した状況を正規乱数で模擬する。両パネルに同じビン境界と共有した軸を用い、分布の位置と幅を比較する。

\begin{lstlisting}[language=Python]
import numpy as np
import matplotlib.pyplot as plt

rng = np.random.default_rng(29)
values1 = rng.normal(loc=0.0, scale=1.0, size=3000)
values2 = rng.normal(loc=0.5, scale=1.2, size=3000)

bin_edges = np.linspace(min(values1.min(), values2.min()),
                        max(values1.max(), values2.max()), 41)
fig, axes = plt.subplots(1, 2, figsize=(10, 4), sharex=True, sharey=True)
axes[0].hist(values1, bins=bin_edges, alpha=0.7, label="sample A")
axes[0].set_title("sample A")
axes[0].set_xlabel("value")
axes[0].set_ylabel("count")

axes[1].hist(values2, bins=bin_edges, alpha=0.7, label="sample B")
axes[1].set_title("sample B")
axes[1].set_xlabel("value")

for ax in axes:
    ax.legend(loc="upper right")
    ax.set_axisbelow(True)

fig.tight_layout()
fig.savefig("subplots_example.pdf")
\end{lstlisting}

この例では、\code{axes} は二つの \code{Axes} を並べた配列であり、\code{for ax in axes:} の部分で凡例位置をそろえる処理を両パネルへ適用している。\code{sharey=True} を付けると図~\ref{fig:subplots-example} のように縦軸の尺度がそろうため、左右の高さを直接比較できる。各パネルの系列が一つだけでも凡例を用意しておけば、後からフィット曲線や参照線を追加しても同じ構成を保てる。

\code{tight\_layout} は余白を自動調整するが、すべての文字やパネルの重なりを解消するとは限らない。タイトルや長いラベルが多い場合には、図のサイズや余白を追加で調整する。複数パネルの配置は、比較対象の対応関係が読み取れるかどうかで判断する。

凡例がデータを隠す場合は、\code{loc="best"} で重なりの少ない位置を選ばせるか、\code{bbox\_to\_anchor} で Axes の外側に配置する。自動調整後も、データ、ラベル、凡例の重なりを保存ファイル上で確認する。

\begin{figure}[tbp]
\centering
\includegraphics[width=0.92\linewidth]{subplots_example.pdf}
\caption{左右に並べたヒストグラムの出力例。両パネルで縦軸の尺度を共有しているため、同じ高さが同じカウントに対応する。}
\label{fig:subplots-example}
\end{figure}

\subsection{エラーバーと両対数軸}

測定値の不確かさを図に示すには、点とともにエラーバーを描く。値が数桁にわたって変化する場合には、対数軸\index{たいすうじく@対数軸|textbf}を用いると分布形状を比較しやすい。ただし、対数軸では 0 以下の値を表示できないため、データ範囲を事前に確認する必要がある。

図~\ref{fig:errorbar-example} にエラーバーと対数プロットの例を示す。

\begin{figure}[htbp]
\centering
\includegraphics[width=0.92\linewidth]{errorbar_example.pdf}
\caption{エラーバーと対数軸の設定例。左は線形軸、右は x 軸と y 軸の両方を対数尺度にした図。}
\label{fig:errorbar-example}
\end{figure}

\begin{lstlisting}[language=Python]
import numpy as np
import matplotlib.pyplot as plt

x_val = np.geomspace(1.0, 100.0, 12)
y_val = 5.0 * x_val ** (-1.2)
y_err = 0.15 * y_val

fig, axes = plt.subplots(1, 2, figsize=(10, 4))

# fmt='o' は点を丸で描き、capsize=3 は端部の横棒を 3 point にする
axes[0].errorbar(x_val, y_val, yerr=y_err, fmt='o', capsize=3, label="data")
axes[0].set_xlabel("x")
axes[0].set_ylabel("y")

axes[1].errorbar(x_val, y_val, yerr=y_err, fmt='o', capsize=3, label="data")
axes[1].set_yscale('log')
axes[1].set_xscale('log')
axes[1].set_xlabel("x")
axes[1].set_ylabel("y")

fig.tight_layout()
fig.savefig("errorbar_example.pdf")
\end{lstlisting}

\subsection{分布図における軸・色・密度の読み取り}

分布図を解釈するときは、軸が表す量と単位を特定したうえで、分布の中心と幅、裾、外れ値の順に確認する。

例えば、図~\ref{fig:distribution-examples} の左パネルでは、ヒストグラムとフィット曲線の残差、および分布の左右対称性を調べる。右パネルでは、短時間側の立ち上がりと長時間側の減衰を同じ時定数で説明できるかを調べる。本文では、図から読み取った特徴と、その特徴が結論を支える論理を記述する。

\begin{figure}[tbp]
\centering
\includegraphics[width=0.92\linewidth]{distribution_examples.pdf}
\caption{Matplotlib による分布図とフィットの例。左はガウシアン型、右は指数分布であり、どちらもヒストグラムの上にフィット曲線を重ね、凡例に分布名と推定パラメータを示している。}
\label{fig:distribution-examples}
\end{figure}

\subsection{画像フォーマットの選択：ラスタとベクタ}\index{matplotlib@Matplotlib!savefig@\code{savefig()}}

解析図の保存形式は、図の要素数、出力先、必要な拡大率に応じて選ぶ。

\begin{itemize}
    \item \textbf{ベクタ形式 (\code{.pdf}, \code{.svg}, \code{.eps})}：線や文字を図形要素として保存するため、拡大しても輪郭がぼやけない。線図や要素数の少ない散布図を論文やレポートへ掲載するときは、\code{savefig("plot.pdf")} のようにベクタ形式で保存する。
    \item \textbf{ラスタ形式 (\code{.png}, \code{.jpg})}：図をピクセルの集合として保存するため、拡大率によっては輪郭に画素が見える。ベクタ形式では各点が個別要素となるので、密集した散布図や各セルを図形として描くヒートマップではファイルサイズと描画負荷が増える。\code{imshow()} の画像は PDF 内にもラスタ画像として埋め込まれる。多数の点やセルだけに \code{rasterized=True} を指定すれば、文字をベクタのまま残すこともできる。図全体を画像にする場合は、\code{savefig("plot.png", dpi=150)} のように必要な解像度を指定して PNG で保存する。
\end{itemize}

\begin{lstlisting}[language=Python]
fig, ax = plt.subplots(figsize=(6, 4))
ax.plot(x, y)
ax.set_xlabel("time [ns]")
ax.set_ylabel("count")

fig.savefig("spectrum.pdf")
fig.savefig("spectrum.png", dpi=150)
\end{lstlisting}

同じ図でも、\code{spectrum.pdf} は線と文字をベクタ要素として保持し、\code{spectrum.png} は指定した画素数のラスタ画像として保存する。線図や要素数の少ないヒストグラムには PDF を用いる。点や格子の数が多く、ベクタ形式ではファイルサイズと描画負荷が増える散布図やヒートマップには、必要な解像度を指定した PNG を用いる。


\section{オブジェクト指向 API による複雑な図の管理}\label{sec:matplotlib-object}

\subsection{Figure, Axes, Axis, Artist の役割}

Matplotlib では、\code{Figure}、\code{Axes}、\code{Axis} が異なる階層を表す。\code{Axes} は複数形に見えるが、Matplotlib では一つの作図領域を表すクラス名である。各オブジェクトの対応を次に示す。

\begin{description}
\item[\code{Figure}] 図全体である。1 つの PDF や PNG、1 つの表示ウィンドウに相当する。\code{savefig}、\code{suptitle}、図全体の凡例や colorbar などは Figure のメソッドまたは構成要素である。
\item[\code{Axes}] データを描く 1 つのパネルである。\code{plot}、\code{hist}、\code{scatter}、\code{set\_xlabel}、\code{legend} などは Axes のメソッドとして呼び出す。
\item[\code{Axis} (\code{x Axis, y Axis})] 目盛り、目盛りラベル、スケールを管理する部品である。通常は \code{ax.set\_xscale("log")} や \code{ax.set\_xlim(...)} のように Axes 経由で触るが、内部には \code{ax.xaxis} と \code{ax.yaxis} がある。
\item[\code{Artist}] 図の中に描かれる要素の総称である。線は \code{Line2D}、文字は \code{Text}、凡例は \code{Legend}、ヒストグラムの棒は \code{Rectangle} という Artist として管理される。代表的な要素を次に示す。
  \begin{itemize}
    \item \code{Title}, \code{Legend}: 図のタイトル、および各系列を示す凡例。
    \item \code{Grid}, \code{Spine}: グリッド線と、座標系を囲む枠線。
    \item \code{Line}, \code{Markers}: 折れ線や散布図の点。
    \item \code{xlabel}, \code{ylabel}: x軸・y軸のラベル。
  \end{itemize}
\end{description}

Figure、Axes、Axis、Artist の階層関係だけを先に図解の構図に合わせて木構造で書くと、次のようになる。

\begin{lstlisting}[numbers=none, xleftmargin=2\zw]
Figure
|-- Axes
|   |-- Title
|   |-- Spine, Grid
|   |-- xaxis, yaxis (x Axis, y Axis / xlabel, ylabel)
|   |-- Line2D, Markers
|   `-- Legend
`-- Figure-level Text / Colorbar / ...
\end{lstlisting}

図~\ref{fig:anatomy-example} に、Figure、Axes、Axis、Artist の位置関係を示す。Figure は一つ以上の Axes を持ち、各 Axes は x 軸と y 軸を表す Axis、および線、文字、凡例などを持つ。Figure、Axes、Axis 自体も Artist の一種であり、四者が互いに独立した分類というわけではない。

\begin{figure}[tbp]
\centering
\includegraphics[width=0.75\linewidth]{anatomy_example.pdf}
\caption{Figure、Axes、Axis、Artist の関係を示した例。1 つの Figure の中に Axes があり、その Axes の中で線、文字、凡例などの Artist が描かれている。}
\label{fig:anatomy-example}
\end{figure}

\noindent この対応を実際の Python オブジェクトの型で示すコードを次に掲げる。

\begin{lstlisting}[language=Python]
import matplotlib.pyplot as plt

x = [0, 1, 2, 3]
y = [1.0, 2.0, 1.5, 3.0]

fig, ax = plt.subplots(figsize=(6, 4))
line, = ax.plot(x, y, label="signal")
ax.set_title("Axes title")
ax.set_xlabel("time")
ax.set_ylabel("signal")
legend = ax.legend()

print(type(fig).__name__)
print(type(ax).__name__)
print(type(ax.xaxis).__name__)
print(type(line).__name__)
print(type(legend).__name__)
print(len(fig.axes))
fig.savefig("line.pdf")
\end{lstlisting}

このコードを実行すると、次の出力が得られる。

\begin{lstlisting}[numbers=none, xleftmargin=2\zw]
Figure
Axes
XAxis
Line2D
Legend
1
\end{lstlisting}

最後の \code{1} は、この Figure の中に Axes が 1 枚だけ入っていることを表す。もし \code{plt.subplots(2, 2)} と書けば、Figure は 1 つのままで、Axes が 4 枚入る。

\code{Artist} は線だけを指す名称ではない。タイトルや軸ラベルは \code{Text}、ヒストグラムの各棒は \code{Rectangle} として管理される。次のコードは、それぞれの返り値の型を表示する。

\begin{lstlisting}[language=Python]
import matplotlib.pyplot as plt

fig, ax = plt.subplots(figsize=(6, 4))
title = ax.set_title("Histogram")
xlabel = ax.set_xlabel("value")
bars = ax.hist([1, 1, 2, 3], bins=[0.5, 1.5, 2.5, 3.5])[2]

print(type(title).__name__)
print(type(xlabel).__name__)
print(type(bars).__name__)
print(type(bars[0]).__name__)
\end{lstlisting}

\noindent このコードの出力は次のとおりである。

\begin{lstlisting}[numbers=none, xleftmargin=2\zw]
Text
Text
BarContainer
Rectangle
\end{lstlisting}

ここで \code{bars} 全体は \code{BarContainer} であり、その中の 1 本 1 本は \code{Rectangle} である。したがって、「ヒストグラムの棒だけ色を変える」といった細かい調整は、最終的には Artist を触る操作だと理解できる。

Figure 側の要素も、コード上の型を見ると役割を区別しやすい。図全体タイトル \code{suptitle} や colorbar は、特定の Axes だけでなく Figure 全体の構成に関わる。次のコードは、これらの型を表示する。

\begin{lstlisting}[language=Python]
import numpy as np
import matplotlib.pyplot as plt

fig, ax = plt.subplots(figsize=(5, 4))
img = ax.imshow(np.arange(4).reshape(2, 2))
suptitle = fig.suptitle("Figure title")
cbar = fig.colorbar(img, ax=ax)

print(type(suptitle).__name__)
print(type(cbar).__name__)
\end{lstlisting}

\noindent このコードでは、次の出力が得られる。

\begin{lstlisting}[numbers=none, xleftmargin=2\zw]
Text
Colorbar
\end{lstlisting}

\code{suptitle} は Figure 全体に付く \code{Text} であり、\code{colorbar} は Figure に追加される補助要素である。複数パネルに共通するタイトルや colorbar は、Figure のメソッドを通じて追加する。

Axis は目盛り、目盛りラベル、表示範囲、尺度を管理する。次の操作は、Axes のメソッドを通じて x Axis の設定を変更する。

\begin{lstlisting}[language=Python]
fig, ax = plt.subplots(figsize=(6, 4))
ax.plot([1, 10, 100], [2, 3, 5], marker="o")
ax.set_xlim(1, 100)
ax.set_xscale("log")
\end{lstlisting}

\code{set\_xlim} は表示範囲を、\code{set\_xscale("log")} は座標値から表示位置への変換を対数尺度へ変更する。多くの設定は、\code{ax.xaxis} を直接操作せず、Axes のメソッドから変更できる。

操作対象とメソッドの対応を次に示す。

\begin{itemize}
\item 図全体を保存する、図全体タイトルを付ける: Figure
\item 線やヒストグラムを描く、軸ラベルや凡例を付ける: Axes
\item 対数軸、目盛り位置、目盛りラベルを変える: Axis 関連。ただし実際の操作は Axes メソッド経由が多い
\item 特定の線だけ色や太さを変える、凡例や文字を細かく調整する: Artist
\end{itemize}

保存には \code{fig.savefig(...)}、凡例には \code{ax.legend()}、対数軸には \code{ax.set\_xscale("log")} を用いるのは、それぞれの操作対象が Figure または Axes だからである。各メソッドの呼び出し元が \code{fig} と \code{ax} のどちらかを追うと、処理対象を判別しやすい。

\subsection{オブジェクト指向 API による描画}

Matplotlib には、現在の作図状態へ描画する \code{pyplot} API と、\code{Figure} と \code{Axes} を明示的に操作するオブジェクト指向 API がある。複数パネルや再利用する作図関数では、描画先を引数として渡せる後者を用いる。

\begin{lstlisting}[language=Python]
import numpy as np
import matplotlib.pyplot as plt

rng = np.random.default_rng(7)
values = rng.normal(loc=0.0, scale=1.2, size=2500)

fig, ax = plt.subplots(figsize=(6, 4))
ax.hist(values, bins=40)
ax.set_xlabel("value")
ax.set_ylabel("count")
fig.savefig("histogram.pdf")
\end{lstlisting}

描画先の \code{ax} を明示すると、同じ処理を複数の Figure や Axes に適用できる。また、関数の引数として Axes を受け取れるため、作図処理を再利用できる。

\code{Figure} と \code{Axes} の区別は、設定の適用範囲にも対応する。図全体のタイトルや余白は Figure、軸ラベルや凡例は各 Axes に設定する。

\subsection{補助関数による作図処理の集約}

同じ体裁の図を反復して作成する場合は、Axes を引数に取る補助関数へ作図処理をまとめる。

\begin{lstlisting}[language=Python]
import numpy as np
import matplotlib.pyplot as plt

def plot_hist(ax, values, label):
    ax.hist(values, bins=40, alpha=0.7, label=label)
    ax.set_xlabel("value")
    ax.set_ylabel("count")
    ax.legend(loc="upper right")

rng = np.random.default_rng(0)
values1 = rng.normal(loc=0.0, scale=1.0, size=2500)
values2 = rng.normal(loc=0.8, scale=1.3, size=2500)

bin_edges = np.linspace(min(values1.min(), values2.min()),
                        max(values1.max(), values2.max()), 41)
fig, axes = plt.subplots(1, 2, figsize=(10, 4), sharex=True, sharey=True)
plot_hist(axes[0], values1, "sample A")
plot_hist(axes[1], values2, "sample B")
fig.savefig("histogram_pair_helper.pdf")
\end{lstlisting}

この補助関数は描画先の Axes を引数として受け取るため、比較図の各パネルや条件ごとの反復描画に同じ処理を適用できる。

\section{ラベル・凡例・色・スタイルの管理}

図の可読性に関わる主な要素は次のとおりである。

\begin{itemize}
\item 量と単位を含む軸ラベル
\item 簡潔なタイトル
\item 適切な余白
\item 統一された色と線幅
\item レポート全体で一貫した表記
\end{itemize}

色、フォント、線幅、軸ラベル、単位、凡例の規則を文書全体で統一すると、同じ表現が同じ意味に対応する。系列や条件の区別には、文書内で一貫した色、マーカー、線種を割り当てる。

\subsection{マーカー、線種、色の指定}

Matplotlib では、点の形\index{marker@marker}、線種\index{linestyle@linestyle}、色をそれぞれ独立に指定できる。対応する主な引数は次の 5 つである。

\begin{lstlisting}[language=Python]
ax.plot(
    x,
    y,
    marker="o",
    markersize=5,
    linestyle="--",
    linewidth=2.0,
    color="C3",
)
\end{lstlisting}

\begin{itemize}
\item \code{marker}: 点の形を決める。
\item \code{markersize}: 点の大きさを point 単位で指定する。\code{scatter()} の \code{s} は point の二乗単位であり、同じ数値でも意味が異なる。
\item \code{linewidth}: 線の太さを point 単位で指定する。
\item \code{linestyle}: 実線、破線、点線などを決める。
\item \code{color}: 線や点の色を決める。
\end{itemize}

図~\ref{fig:style-reference-example} は、マーカー、線種、既定の色循環の指定と描画結果を並べたものである。

\begin{figure}[tbp]
\centering
\includegraphics[width=0.92\linewidth]{style_reference_example.pdf}
\caption{マーカー、線種、既定色の代表例。各指定と描画結果の対応を示す。}
\label{fig:style-reference-example}
\end{figure}

図~\ref{fig:style-reference-example} では、マーカー、線種、色の各指定と描画結果を比較できる。個々の指定と用途は、マーカーを表~\ref{tab:matplotlib-markers}、線種を表~\ref{tab:matplotlib-linestyles}、色を表~\ref{tab:matplotlib-colors} にまとめた。系列の識別を色だけに依存させず、マーカーまたは線種も併用すれば、白黒印刷でも系列間の差を保持できる。

\begin{table}[tbp]
\centering
\caption{marker の指定例}
\label{tab:matplotlib-markers}
\begin{tabular}{lll}
\toprule
指定 & 形状 & 用途の例 \\
\midrule
\code{"o"} & 丸 & 測定点、既定の散布図 \\
\code{"s"} & 四角 & 条件の違いを強調したい点 \\
\code{"\textasciicircum"} & 上向き三角 & 系列ごとの差別化 \\
\code{"D"} & ひし形 & 参照データや補助系列 \\
\code{"x"} & バツ & 強調した外れ値や選択点 \\
\bottomrule
\end{tabular}
\end{table}

\begin{table}[tbp]
\centering
\caption{line style の指定例}
\label{tab:matplotlib-linestyles}
\begin{tabular}{lll}
\toprule
指定 & 線種 & 用途の例 \\
\midrule
\code{"-"} & 実線 & 主結果、連続関数 \\
\code{"--"} & 破線 & フィット曲線、比較対象 \\
\code{":"} & 点線 & 補助線、閾値 \\
\code{"-."} & 一点鎖線 & 参照モデル、別条件 \\
\bottomrule
\end{tabular}
\end{table}

\begin{table}[tbp]
\centering
\caption{色指定の基本}
\label{tab:matplotlib-colors}
\begin{tabular}{lll}
\toprule
指定 & 例 & 備考 \\
\midrule
\code{"C0"} など & \code{"C3"} & 既定の color cycle を順に使う \\
名前付き色 & \code{"black"}, \code{"tab:blue"} & 手早く書ける \\
16 進数 & \code{"\#d62728"} & RGB の色を明示的に固定する場合 \\
\bottomrule
\end{tabular}
\end{table}

色だけによる系列の区別は、白黒印刷や色覚特性によって失われる場合がある。系列が 2 つ以上ある場合は、色に加えて marker または line style を変え、色以外からも系列を識別できるようにする。

\subsection{\code{rcParams} による描画設定の統一}

Matplotlib では、図のフォントや線幅などの既定値を \code{rcParams}\index{rcparams@\code{rcParams}} という辞書オブジェクトで管理する。複数の図に共通する設定は、作図スクリプトの冒頭でこの辞書へ指定する。

\begin{lstlisting}[language=Python]
import matplotlib as mpl

mpl.rcParams.update(
    {
        "font.family": "serif",
        "mathtext.fontset": "cm",
        "font.size": 12,
        "axes.grid": True,
        "grid.linestyle": "--",
        "xtick.direction": "in",
        "ytick.direction": "in",
        "axes.linewidth": 1.2,
        "legend.frameon": False,
    }
)
\end{lstlisting}

\begin{itemize}
\item \code{font.family}: 既定フォント族を指定する。本文と図で同じフォント族を使う場合に設定する。
\item \code{font.size}: 文字の基準サイズを指定する。出力媒体と図の寸法に応じて設定する。
\item \code{axes.grid}, \code{grid.linestyle}: グリッドを既定で表示するか、その線種をどうするかを決める。
\item \code{xtick.direction}, \code{ytick.direction}: 目盛り線を内向きにするか外向きにするかを決める。
\item \code{axes.linewidth}: Axes を囲む枠線の太さを指定する。
\item \code{legend.frameon}: 凡例の枠を出すかどうかを決める。
\end{itemize}

\code{rcParams} で設定した値は、その後に作成する図の既定値になる。図ごとに異なる項目は、\code{ax.set\_*} などのメソッドで上書きする。

\subsection{\code{cmap} の選び方}

ヒートマップや散布図で色を値へ対応付けるときは、\code{cmap=} でカラーマップを選ぶ。どの色を選んでもよいわけではなく、データの意味に応じて使い分ける必要がある。代表例を表~\ref{tab:matplotlib-colormaps} に示す。

\begin{table}[tbp]
\centering
\caption{カラーマップの分類と指定例}
\label{tab:matplotlib-colormaps}
\begin{tabular}{lll}
\toprule
種類 & 例 & 対応させる量 \\
\midrule
Sequential & \code{viridis}, \code{plasma}, \code{cividis} & 0 から大へ単調に増える量 \\
Diverging & \code{coolwarm}, \code{RdBu\_r} & 0 や基準値をはさんで符号が変わる量 \\
Grayscale & \code{Greys} & 印刷や下地としての濃淡表示 \\
Categorical & \code{tab10} & 離散的な系列の識別 \\
\bottomrule
\end{tabular}
\end{table}

カウント数や強度のように順序と大小を持つ量には、\code{viridis} などの sequential map を対応させる。正負のずれや差分のように基準値の両側を区別する量には、diverging map を対応させる。値の大小を持たない系列の分類には categorical map を用いる。

\subsection{設定ファイル (.mplstyle) による描画設定の共通化}\index{matplotlib@Matplotlib!.mplstyle@\code{.mplstyle}|textbf}

解析テーマごとに作図スクリプトが増えると、各 Python ファイルの先頭へ同じ設定辞書をコピーする方法は保守しにくい。Matplotlib では、共通の描画設定を \code{.mplstyle} 形式のスタイル設定ファイルへ分離し、まとめて読み込める。

作業ディレクトリに置く \code{my\_style.mplstyle} の記述例を次に示す。

\begin{lstlisting}[numbers=none, xleftmargin=2\zw]
font.family: serif
mathtext.fontset: cm
font.size: 12
axes.grid: True
grid.linestyle: --
xtick.direction: in
ytick.direction: in
axes.linewidth: 1.2
legend.frameon: False
\end{lstlisting}

各解析スクリプトの先頭でこの設定ファイルを読み込むと、以降の図に共通の設定が適用される。

\begin{lstlisting}[language=Python]
import matplotlib.pyplot as plt

# スクリプトの冒頭で自分のスタイルシートを読み込む
plt.style.use("./my_style.mplstyle")

# 以降の描画はすべて統一されたスタイル（Serif フォント等）が反映される
fig, ax = plt.subplots(figsize=(6, 4))
\end{lstlisting}

設定をファイルへ分離すると、解析スクリプトごとのフォントやグリッドの不統一を防げる。学会発表用に文字を大きくした別の設定へも容易に切り替えられる。

\subsection{共通スタイルの部分的な上書き}

共通スタイルを直接変更すると、そのスタイルを読むすべての図へ変更が及ぶ。発表用に文字だけを拡大する場合や、下書きだけグリッドを消す場合は、基準スタイルを読み込んだ後に必要な項目だけを上書きする。

\begin{lstlisting}[numbers=none, xleftmargin=2\zw]
# presentation_style.mplstyle
font.size: 16
axes.grid: False
legend.fontsize: 11
\end{lstlisting}

\begin{lstlisting}[language=Python]
import matplotlib as mpl
import matplotlib.pyplot as plt

plt.style.use(["./my_style.mplstyle", "./presentation_style.mplstyle"])

# スタイルファイルを増やしたくないなら、その場で個別に上書きしてもよい
mpl.rcParams["axes.titlesize"] = 14

fig, ax = plt.subplots(figsize=(6, 4))
ax.grid(False)  # この Axes だけグリッドを消す
\end{lstlisting}

二つの \code{.mplstyle} を配列で渡すと、後ろに書いたファイルの設定が優先される。したがって、基準となる \code{my\_style.mplstyle} を保ったまま、文字サイズやグリッドの有無だけを差し替えられる。さらに、\code{mpl.rcParams} の変更はその後に作成する図へ、\code{ax.grid(False)} のような Axes メソッドは指定したパネルへ適用される。スタイルファイル、スクリプト内の既定値、個別の Axes の順に、適用範囲が狭くなる。


\section{図中の数式表記}

\subsection{Matplotlib による数式表記}

Matplotlib では、軸ラベルや凡例に数式風の表記を書ける。既定では LaTeX そのものではなく、TeX に似た \code{MathText}\index{mathtext@MathText} という仕組みが使われる。

\begin{lstlisting}[language=Python]
plt.xlabel(r"$\Delta t\,[\mathrm{ns}]$")
plt.ylabel(r"$N$")
\end{lstlisting}

単位を数式の中へ入れる場合は、上の例のように \code{\textbackslash{}mathrm} を使って立体で表す。図中の記号、字体、単位は、本文での定義と一致させる。

MathText は、Matplotlib 内部で数式の一部を組版する機能であり、LaTeX の処理系そのものではない。分数、上付き・下付き、ギリシャ文字、\code{\textbackslash{}mathrm} などを扱える一方、任意の LaTeX マクロやパッケージ依存の記法には対応しない。本文と同じ LaTeX 処理系による組版が必要な場合は、外部依存関係を記録したうえで \code{text.usetex=True} を検討する。

本文で $\tau$ を時定数として定義したなら、図の凡例とキャプションでも同じ記号を用いる。解析中の量と図中の表記を一対一に対応させる。


\section{\code{pyplot} と \code{Axes}・カテゴリ軸・画像保存の注意点}

\subsection{\code{plt} と \code{ax} の混在による参照先の混乱}

\code{plt.xlabel(...)} のような状態依存 API は、現在選択されている Axes を変更する。複数の Axes や Figure を扱うコードでは、\code{fig, ax = plt.subplots(...)} で対象を受け取り、\code{ax.set\_xlabel(...)} のようなメソッドで変更先を明示する。

\subsection{文字列データによるカテゴリ軸への変換}

数字に見える列でも、実際には文字列として読まれていると、Matplotlib は数値軸ではなくカテゴリ軸として解釈する。すると、\code{1}, \code{2}, \code{10}, \code{20} が数値間隔ではなく、単なる 4 つのラベルとして等間隔に並ぶ。

\begin{lstlisting}[language=Python]
import matplotlib.pyplot as plt

x_str = ["1", "2", "10", "20"]
y = [1.0, 4.0, 2.0, 5.0]
fig, ax = plt.subplots(figsize=(6, 4))
ax.plot(x_str, y, marker="o")
\end{lstlisting}

図~\ref{fig:string-axis-example} の左では文字列がカテゴリ軸として扱われ、右では数値への変換によって数値軸になっている。CSV の読み込み後に軸間隔が数値に比例しない場合は、作図に渡した配列の \code{dtype} を確認する。

\begin{figure}[tbp]
\centering
\includegraphics[width=0.92\linewidth]{string_axis_example.pdf}
\caption{文字列軸と数値軸の違い。左は \code{"1"}, \code{"2"}, \code{"10"}, \code{"20"} をカテゴリとして並べており、右は数値として扱っている。}
\label{fig:string-axis-example}
\end{figure}

\subsection{\code{savefig} と余白調整の順序}

長いラベルや図外の凡例は、保存領域からはみ出す場合がある。\code{tight\_layout()}、\code{constrained\_layout=True}、または \code{bbox\_inches="tight"} の効果を保存したファイルで確認する。再現可能なスクリプトでは、描画、余白調整、保存の順序を固定する。対数軸へ渡す値は正でなければならない。

\subsection{公式資料における作図例の検索}

ここまでの各節では、基本的な描画、レイアウト、オブジェクトの役割、保存方法を説明した。実際の作図では、目的に近い例を公式資料から探すと、利用できる図種や引数を把握しやすい。特に次の 4 つは参照頻度が高い。

\begin{itemize}
\item Quick start~\cite{matplotlib_quick_start}
\item Plot types~\cite{matplotlib_plot_types}
\item Gallery~\cite{matplotlib_gallery}
\item Cheatsheets~\cite{matplotlib_cheatsheets}
\end{itemize}

\code{Plot types} は図種ごとの選択肢、\code{Gallery} は完成例と実装、\code{Cheatsheets} は主要な引数と書式、\code{Quick start} は \code{Figure} と \code{Axes} の基本構造を示す。目的に近い公式例を選び、入力データ、軸、単位、保存形式を解析条件に合わせて変更する。

\section{章末課題：地震カタログと発生間隔の可視化}

USGS データと第~\ref{chap:numpy}~章で作成した \(\Delta t\) 配列がない場合は、次の小規模データを共通入力に用いよ。課題では画面表示だけで終わらせず、指定したファイルを保存し、保存後に \code{Path(...).exists()} で存在を確認せよ。

\begin{lstlisting}[language=Python]
import numpy as np

longitude = np.array([140.0, 141.2, 139.5, 142.0, 138.7])
latitude = np.array([35.0, 36.1, 34.8, 37.0, 33.9])
magnitude = np.array([3.2, 4.1, 2.8, 4.8, 3.7])
dt = np.array([1.0, 4.0, 12.0, 60.0, 180.0, 900.0, 3600.0])
\end{lstlisting}

\begin{enumerate}
\item \textbf{全球地震分布の散布図}。
USGS の CSV から \code{longitude}、\code{latitude}、\code{mag} を読み、横軸を経度、縦軸を緯度とする散布図を作れ。点の色でマグニチュードを表し、軸ラベル、カラーバー、タイトルを付けたうえで \code{figures/usgs\_global\_scatter.pdf} へ保存せよ。
\code{fig, ax = plt.subplots()} の返り値、\code{scatter = ax.scatter(...)} の返り値、\code{fig.colorbar(scatter, ax=ax)} の第 1 引数と \code{ax} 引数の役割を説明せよ。\code{scatter.get\_array()} がマグニチュード配列と一致し、\code{len(fig.axes)} がカラーバーを含めて 2 になることを確認せよ。

\item \textbf{発生間隔ヒストグラムの対数表示}。
第~\ref{chap:numpy}~章で作った \(\Delta t\) 配列を用い、発生間隔ヒストグラムを描け。縦軸を対数表示とし、横軸を線形尺度とした場合と対数尺度とした場合について、短い発生間隔の分解能、長時間側の裾の範囲、空のビンの数を比較せよ。図中にはビンの意味が分かる軸ラベルと単位を明記せよ。
線形ビンの図と \code{np.logspace()} で作った対数ビンの図を 2 枚の Axes に並べ、\code{figures/interval\_comparison.pdf} へ保存せよ。\code{ax.set\_yscale("log")} と \code{ax.set\_xscale("log")} がデータ値そのものを変換するのではなく、Axis の表示尺度を変えることを説明せよ。0 以下の値を対数軸へ置けない理由も述べよ。

\item \textbf{発展課題}。
通常の散布図では、点がどの海溝や列島に沿って並んでいるかは読み取りにくい。海岸線付きの地図へ重ねたい場合は、Matplotlib だけではなく地図投影を扱う \code{cartopy} が必要になる。\code{cartopy} を導入し、同じ地震データを海岸線付きの地図へ重ねよ。必要な出発点は、
\begin{itemize}
\item \code{projection=ccrs.PlateCarree()} による地図 Axes の作成
\item \code{scatter} 側への \code{transform=ccrs.PlateCarree()} の指定
\item \code{ax.coastlines()} による海岸線の描画
\item 全球図における \code{ax.set\_global()} の使用
\item 日本近海図における \code{ax.set\_extent([122, 154, 20, 50], crs=ccrs.PlateCarree())} の使用
\end{itemize}
である。これらの使い方を検索し、地図付きの散布図を作成せよ。
最小コードは、\code{import cartopy.crs as ccrs}、\code{fig, ax = plt.subplots(subplot\_kw=\{"projection": ccrs.PlateCarree()\})}、\code{ax.set\_global()}、\code{ax.coastlines()}、\code{ax.scatter(longitude, latitude, c=magnitude, transform=ccrs.PlateCarree())} の順に組み立てられる。\code{projection} は地図を描く Axes の座標系、\code{transform} は入力データの座標系を指定する。\code{figures/usgs\_map.pdf} を保存し、両者を同じ値にしていても役割が異なることを説明せよ。Cartopy が未導入なら、実行した導入コマンドと、導入前に出た例外名を記録せよ。
\end{enumerate}

可視化によって分布の形、外れ値、残差の傾向を確認できると、次に検討すべき数値計算を選びやすくなる。次章では、図から得た仮説を定量化するために、SciPy の統計分布、フィット、補間、信号処理を扱う。
